% Lösungen Prozentrechnung (zusammenbau v0.4, Heft)
\einheitenkopf{Prozentrechnung}

\erg{23}{a) $50\,\%$ \quad b) $37\,\%$}
\erg{24}{a) $5$ von $100$ Paketen kommen zu spät. ($5\,\% = \frac{5}{100}$; „jedes 5.“ wären $20\,\%$) \quad b) $8$ von $100$ Gästen bestellen Tee. ($8\,\% = \frac{8}{100}$; „jeder 8.“ wären $12{,}5\,\%$) \quad c) $25\,\%$ ($\frac{1}{4} = \frac{25}{100}$) \quad d) $5\,\%$ ($\frac{1}{20} = \frac{5}{100}$)}
\erg{25}{a) Aussage 2 ist falsch: $200 : 600 \approx 0{,}33$, also ein Drittel. Richtig: Ein Drittel der Kinder fährt mit dem Rad. (Aussage 1 stimmt: $150 : 600 = 0{,}25$.) \quad b) Aussage 1 ist falsch: $600 : 2\,400 = 0{,}25$, also ein Viertel. Richtig: Jedes vierte Mitglied ist ein Kind. (Jedes fünfte, $480 : 2\,400 = 0{,}2$, gilt für die Jugendlichen; Aussage 2 stimmt: $1\,200 : 2\,400 = 0{,}5$.) \quad c) Aussage 1 ist falsch: auf dem Land sind es $\frac{1}{10} = 10\,\%$. Richtig: Ein Zehntel der Befragten auf dem Land geht in die Bibliothek (ein Viertel, $\frac{1}{4} = 25\,\%$, gilt für die Stadt). Aussage 2 stimmt: $55\,\% > 50\,\%$. \quad d) Aussage 2 ist falsch: $\frac{1}{20} = 5\,\%$. Richtig: Jeder zwanzigste Grundschüler frühstückt nie zu Hause (jeder zehnte gilt für die Oberschule). Aussage 1 stimmt: $\frac{3}{5} = 60\,\%$. \quad e) Nein: $27{,}5\,\% + 22{,}0\,\% = 49{,}5\,\%$, das ist weniger als die Hälfte. \quad f) Ja: $45{,}3\,\% + 21{,}8\,\% + 4{,}3\,\% = 71{,}4\,\%$ (oder $100\,\% - 28{,}6\,\%$), das ist mehr als $70\,\%$.}
\erg{26}{a) Biogas: $100 - 38 - 24 - 9 - 3 = 26$, also $26\,\%$ \quad b) Siedlung: $100 - 34 - 41 - 4 - 2 = 19$, also $19\,\%$}
\erg{27}{a) $20\,\%$ \quad b) $17\,\%$ \quad c) $\frac{19}{50} = \frac{38}{100} = 38\,\%$}
\erg{28}{a) Ganzes $24$ Lose; $3 : 24 = 0{,}125 = 12{,}5\,\%$ \quad b) Ganzes $15$; $4 : 15 \approx 0{,}267 \approx 26{,}7\,\%$ \quad c) $57\,000 : 480\,000 = 0{,}11875 \approx 11{,}9\,\%$ \quad d) $136 : 380 \approx 0{,}358 \approx 35{,}8\,\%$ \quad e) abgelehnt $64 - 47 = 17$; $17 : 64 \approx 0{,}266 \approx 26{,}6\,\%$ \quad f) verspätet $92 - 79 = 13$; $13 : 92 \approx 0{,}141 \approx 14{,}1\,\%$ \quad g) $1\,080 : 40\,000 = 0{,}027 = 2{,}7\,\%$ \quad h) $175 : 12\,500 = 0{,}014 = 1{,}4\,\%$ \quad i) $6$ von $13$ sind größer; $6 : 13 \approx 0{,}462 \approx 46{,}2\,\%$ \quad j) $5$ von $9$ warteten länger; $5 : 9 \approx 0{,}556 \approx 55{,}6\,\%$}
\erg{29}{a) $40\,€$ \quad b) $1\,\% = 4$ kg; $6\,\% = 24$ kg}
\erg{30}{a) $72\,€$ ($480 \cdot 0{,}15 = 72$; $408\,€$ ist der Restpreis) \quad b) $195\,€$ ($650 \cdot 0{,}3 = 195$; $455\,€$ ist der Restpreis) \quad c) $240 \cdot 0{,}15 = 36\,€$ (nicht der Restpreis $204\,€$) \quad d) $180 \cdot 0{,}35 = 63\,€$ (nicht der Restpreis $117\,€$) \quad e) $60 \cdot 0{,}17 = 10{,}20\,€$ \quad f) $40 \cdot 0{,}19 = 7{,}60\,€$ \quad g) $0{,}4 \cdot 85 = 34\,€$ (nicht $0{,}04$) \quad h) $0{,}6 \cdot 45 = 27\,€$ (nicht $0{,}06$) \quad i) $52\,\% - 35\,\% = 17$ Prozentpunkte; $0{,}17 \cdot 800 = 136$ Haushalte (nicht $17$) \quad j) $44\,\% - 26\,\% = 18$ Prozentpunkte; $0{,}18 \cdot 650 = 117$ Schüler (nicht $18$)}
\erg{31}{a) $20$ Kästchen \quad b) $2 \cdot 35 = 70$ kg \quad c) $1\,\% = 4$ kg; $100\,\% = 400$ kg}
\erg{32}{a) $4 \cdot 4 = 16\,€$ (nicht $25\,\%$ von $4\,€$) \quad b) $1\,\% = 0{,}50\,€$; $100\,\% = 50\,€$ (nicht $30\,\%$ von $15\,€$) \quad c) $1\,\% = 150\,€$; $100\,\% = 15\,000\,€$ (nicht $40\,\%$ von $6\,000\,€$) \quad d) $5 \cdot 450 = 2\,250$ kg (nicht $20\,\%$ von $450$ kg)}
\erg{33}{a) auf $80\,\%$ gesenkt \quad b) $10\,\% = 8\,€$; $80 + 8 = 88\,€$ \quad c) $45 \cdot 1{,}2 = 54\,€$}
\erg{34}{a) $1{,}29 - 1{,}15 = 0{,}14$; $0{,}14 : 1{,}15 \approx 0{,}122 \approx 12{,}2\,\%$ \quad b) $1{,}90 - 1{,}70 = 0{,}20$; $0{,}20 : 1{,}70 \approx 0{,}118 \approx 11{,}8\,\%$ \quad c) $1\,340 - 1\,120 = 220$; $220 : 1\,340 \approx 0{,}164 \approx 16{,}4\,\%$ \quad d) $860 - 780 = 80$; $80 : 860 \approx 0{,}093 \approx 9{,}3\,\%$ \quad e) $57\,600 - 48\,000 = 9\,600$; $9\,600 : 48\,000 = 0{,}2 = 20\,\%$ (nicht durch den neuen Wert) \quad f) $72\,000 - 64\,000 = 8\,000$; $8\,000 : 64\,000 = 0{,}125 = 12{,}5\,\%$ (nicht durch den neuen Wert) \quad g) $2{,}40 \cdot 1{,}25 = 3{,}00\,€$ (nicht nur $0{,}60\,€$) \quad h) $65 \cdot 1{,}12 = 72{,}80\,€$ (nicht nur $7{,}80\,€$) \quad i) voller Preis $2 \cdot 14 + 2 \cdot 9 = 46\,€$; $46 \cdot 0{,}85 = 39{,}10\,€$ (nicht der Rabatt $6{,}90\,€$) \quad j) voller Preis $3 \cdot 22 + 13 = 79\,€$; $79 \cdot 0{,}75 = 59{,}25\,€$ (nicht der Rabatt $19{,}75\,€$) \quad k) Lücken: $100$ m und $5$ m; $15 : 240 = 0{,}0625 \approx 6{,}3\,\% > 5\,\%$ – Herr Lenz hat recht. \quad l) $15$ m; $1{,}1 : 8 = 0{,}1375 \approx 13{,}8\,\% < 15\,\%$ – Frau Roth hat nicht recht, die Zufahrt ist zulässig.}
